% ======================================================================
%  KAVRAM: Embedding  (Tier 1)
% ======================================================================

\concept{Knowledge · Bilgi}{Embedding}{Anlamı koordinata çevirmek: benzer anlam, yakın konum}

% ----------------------------------------------------------------------
\stage{1}{Problem}
% ----------------------------------------------------------------------
Bilgisayar için ``araba'' ve ``otomobil'' tamamen alakasız iki karakter
dizisidir --- ortak harf bile yok. Oysa anlamca aynı şey. Metni harf harf
karşılaştırdığın sürece, eş anlamlıları, yakın kavramları, aynı soruyu farklı
kelimelerle soran kullanıcıyı asla yakalayamazsın.

\keyline{Sorun: anlam, yazımın içinde değildir. Anlamı aramak için metni anlamın yaşadığı bir uzaya taşımamız gerekir.}

Klasik anahtar kelime araması (\en{keyword search}) tam da burada çöker:
kullanıcı ``yağmurluk'' arar, kataloğunda ``su geçirmez mont'' yazar, ve hiçbir
sonuç dönmez --- aynı şey olmasına rağmen.

% ----------------------------------------------------------------------
\stage{2}{Neden Var}
% ----------------------------------------------------------------------
Anlama göre arama, kümeleme ve karşılaştırma yapabilmek için metni \emph{sayısal}
ve \emph{anlam-koruyan} bir biçime çevirmemiz gerekir. Embedding tam bunu yapar:
benzer anlamlı metinleri uzayda birbirine yakın noktalara yerleştirir.

Bu sayede ``benzerlik'' artık bir geometri sorusuna dönüşür: iki metin ne kadar
yakınsa, anlamca o kadar ilişkilidir. Retrieval, öneri, kümeleme, deduplikasyon ---
hepsi bu tek fikrin üzerine kurulur.

% ----------------------------------------------------------------------
\stage{3}{Tarihçe}
% ----------------------------------------------------------------------
\begin{timeline}
\yr{2000'ler} TF-IDF / bag-of-words: yazımı yakalar, anlamı değil \tlsep \yr{2013} word2vec / GloVe: ``kral − erkek + kadın ≈ kraliçe'' \tlsep \yr{2018} BERT: bağlama duyarlı temsiller \tlsep \yr{2022–26} cümle/doküman embedding API'leri → RAG'in temeli
\end{timeline}

% ----------------------------------------------------------------------
\stage{4}{Tanım}
% ----------------------------------------------------------------------
\begin{problembox}
\textbf{Embedding:} Bir metin parçasının anlamını temsil eden, sabit uzunlukta
bir sayı vektörü --- öyle ki anlamsal benzerlik $\approx$ geometrik yakınlık
(genellikle \term{cosine similarity} ile ölçülür).
\end{problembox}

Örneğin bir model her metni 1536 boyutlu bir vektöre çevirebilir. ``Araba'' ve
``otomobil'' bu uzayda neredeyse üst üste düşerken, ``araba'' ve ``muz''
birbirinden uzaktır.

% ----------------------------------------------------------------------
\stage{5}{Zihinsel Model}
% ----------------------------------------------------------------------
\begin{mentalbox}
Devasa bir \textbf{kütüphane} düşün --- ama kitaplar alfabetik değil,
\emph{konuya} göre yerleştirilmiş. Yemek kitapları bir köşede, astrofizik başka
bir köşede; birbirine yakın fikirler fiziksel olarak da yan yana.

\smallskip
Embedding, bir metnin bu ``anlam haritasındaki koordinatıdır''. Arama yapmak
artık raf numarası eşleştirmek değil, \emph{haritada en yakın komşulara bakmaktır}.
\end{mentalbox}

% ----------------------------------------------------------------------
\stage{6}{Mimari Diyagram}
% ----------------------------------------------------------------------
Solda metni vektöre çeviren boru hattı; sağda anlamın 2 boyuta indirgenmiş hâli ---
benzer kavramlar kümeleniyor, sorgu en yakın komşularını buluyor.

\vspace{4pt}
\begin{center}
\begin{tikzpicture}[font=\sffamily\scriptsize]
  % --- Sol: pipeline ---
  \node[draw=softgray, fill=white, rounded corners=8pt, minimum height=8mm,
        minimum width=20mm, align=center] (txt) {``su geçirmez\\mont''};
  \node[draw=accent, fill=wash, rounded corners=2pt, below=6mm of txt,
        minimum height=10mm, minimum width=24mm, align=center] (emb)
        {Embedding\\modeli};
  \node[draw=accentdk, fill=accent!12, rounded corners=2pt, below=6mm of emb,
        minimum height=8mm, minimum width=24mm, align=center, font=\ttfamily\tiny]
        (vec) {[0.21, -0.34,\\0.88, ...]};
  \draw[-{Stealth[length=2mm]}, accent, line width=0.8pt] (txt) -- (emb);
  \draw[-{Stealth[length=2mm]}, accent, line width=0.8pt] (emb) -- (vec);

  % --- Sağ: 2B anlam uzayı ---
  \begin{scope}[shift={(52mm,-20mm)}]
    \draw[hairline, line width=0.8pt, -{Stealth[length=1.6mm]}] (0,0) -- (52mm,0);
    \draw[hairline, line width=0.8pt, -{Stealth[length=1.6mm]}] (0,0) -- (0,40mm);
    \node[anchor=south west, font=\sffamily\bfseries\tiny, text=softgray]
          at (0,40mm) {ANLAM UZAYI \,(2B'ye indirgenmiş)};

    % Giyim kümesi (yayılmış; etiketler 4 farklı yönde)
    \draw[accent, dashed, line width=0.5pt] (11mm,30mm) circle (9.5mm);
    \fill[accent] (17mm,35mm) circle (1.1mm) node[right=1mm,font=\tiny,text=ink]{yağmurluk};
    \fill[accent] (9mm,28mm)  circle (1.1mm) node[below left=0mm and -0.5mm,font=\tiny,text=ink]{mont};
    \fill[accent] (5mm,24mm)  circle (1.1mm) node[left=1mm,font=\tiny,text=ink]{ceket};

    % Taşıt kümesi
    \fill[good] (40mm,10mm) circle (1.1mm) node[right=1mm,font=\tiny,text=ink]{araba};
    \fill[good] (44mm,13mm) circle (1.1mm) node[right=1mm,font=\tiny,text=ink]{otomobil};
    \fill[good] (38mm,7mm) circle (1.1mm) node[below=0.5mm,font=\tiny,text=ink]{kamyon};

    % Sorgu (kırmızı, etiket yukarıda)
    \fill[warn] (14mm,32mm) circle (1.3mm);
    \node[font=\sffamily\tiny, text=warn, above=0.5mm of {(14mm,32mm)}] {sorgu};
  \end{scope}
\end{tikzpicture}
\end{center}

\keyline{``Yağmurluk'' sorgusu, harf benzerliğine değil anlam yakınlığına göre ``su geçirmez mont''un yanına düşer.}

% ----------------------------------------------------------------------
\stage{7}{Karar Çerçevesi}
% ----------------------------------------------------------------------
Hangi embedding modeli, ne zaman yeniden embed? Yanlış seçim tüm retrieval'ı
sessizce bozar.

\vspace{2pt}
\begin{center}
\begin{tikzpicture}[
   font=\sffamily\scriptsize, node distance=5mm and 8mm,
   q/.style={draw=accent, line width=0.8pt, fill=wash, diamond, aspect=2.4,
             align=center, inner sep=1pt, minimum width=30mm},
   leaf/.style={draw=accentdk, line width=0.8pt, rounded corners=2pt,
                fill=accent!12, align=center, inner sep=4pt, minimum width=24mm},
   yes/.style={good, font=\sffamily\bfseries\scriptsize},
   no/.style={warn, font=\sffamily\bfseries\scriptsize},
   ed/.style={-{Stealth[length=2mm]}, softgray, line width=0.7pt}]

  \node[q] (q1) {Çok dilli içerik\\var mı?};
  \node[leaf, below left=14mm and -4mm of q1] (multi) {\textbf{Çok dilli model}};
  \node[q, below right=14mm and -4mm of q1] (q2) {Alan çok özel mi?\\(tıp, hukuk)};
  \node[leaf, below left=13mm and 2mm of q2, minimum width=22mm] (domain) {\textbf{Alan modeli}\\/ fine-tune};
  \node[leaf, below right=13mm and 2mm of q2, minimum width=22mm] (general) {\textbf{Genel model}\\(varsayılan)};

  \draw[ed] (q1) -- node[yes,left]{Evet} (multi);
  \draw[ed] (q1) -- node[no,right]{Hayır} (q2);
  \draw[ed] (q2) -- node[yes,left]{Evet} (domain);
  \draw[ed] (q2) -- node[no,right]{Hayır} (general);
\end{tikzpicture}
\end{center}

\keyline{Altın kural: İndeks ve sorgu \emph{aynı} embedding modeliyle üretilmeli. Model değişirse tüm indeks yeniden embed edilmeli.}

% ----------------------------------------------------------------------
\stage{8}{Üretim Senaryosu}
% ----------------------------------------------------------------------
\begin{prodbox}[Üretim Senaryosu · Anlamsal Ürün Araması]
\textbf{Bağlam.} Bir e-ticaret sitesinde arama anahtar kelimeye dayalı.
``Yağmurluk'' arayan kullanıcı, ürünler ``su geçirmez mont'' diye etiketlendiği
için boş sonuç görüyor; satış kaçıyor.

\smallskip
\textbf{Çözüm.} Tüm ürün açıklamaları bir embedding modeliyle vektöre çevrilip
vektör DB'ye yazılır. Arama sorgusu da aynı modelle embed edilir; en yakın
komşular getirilir.

\smallskip
\textbf{Sonuç.} ``Yağmurluk'', ``su geçirmez mont''u bulur; ``hediye'' aramak
ilgili kategorileri yüzeye çıkarır. Arama artık yazımı değil \textbf{niyeti}
yakalar.
\end{prodbox}

% ----------------------------------------------------------------------
\stage{9}{Anti-pattern}
% ----------------------------------------------------------------------
\begin{antibox}[Anti-pattern · ``Karışık Uzay'']
İndeksi bir modelle, sorguyu başka bir modelle (veya model güncellendikten
sonra) embed etmek; ya da kocaman bir metni tek vektöre sıkıştırıp anlamı
ortalama bir bulamaca çevirmek. İki vektör farklı uzaylardaysa ``yakınlık''
anlamsızdır.

\smallskip
\textbf{Belirti:} ``Retrieval alakasız sonuçlar getiriyor, neden bilmiyoruz.''\\
\textbf{Kök neden:} İndeks ve sorgu aynı anlam uzayında değil; veya chunk çok büyük.\\
\textbf{Düzeltme:} Tek model sabitle, sürümü kaydet; makul boyutta chunk'la; model değişiminde yeniden embed et.
\end{antibox}

% ----------------------------------------------------------------------
\stage{10}{Alternatifler \& Ödünleşimler}
% ----------------------------------------------------------------------
\vspace{2pt}
\begin{tabularx}{\linewidth}{@{}>{\raggedright\arraybackslash}p{30mm} L L@{}}
\arrayrulecolor{hairline}\toprule
{\sffamily\bfseries\small Yaklaşım} &
{\sffamily\bfseries\small Güçlü olduğu yer} &
{\sffamily\bfseries\small Ödünleşim} \\
\midrule
\textbf{Lexical (BM25)} & Tam terim, kod, isim, nadir kelime &
Eş anlam/parafrazı kaçırır \\
\addlinespace[3pt]
\textbf{Dense embedding} & Anlam, eş anlam, niyet &
Tam terimde zayıf; model bağımlı \\
\addlinespace[3pt]
\textbf{Hybrid} & İkisinin gücünü birleştirir &
Daha karmaşık; ayar gerektirir \\
\bottomrule
\end{tabularx}

\vspace{4pt}
Production'da en sağlam retrieval genelde \textbf{hybrid}'dir: anlamsal yakınlık +
tam terim eşleşmesi.

% ----------------------------------------------------------------------
\stage{11}{İlişkili Kavramlar}
% ----------------------------------------------------------------------
\begin{itemize}
  \item \term{Vector Database} --- embedding'leri saklar ve hızlı komşu araması yapar.
  \item \term{Similarity Search} --- en yakın vektörleri bulma işlemi.
  \item \term{Cosine Similarity} --- iki vektörün anlamsal yakınlık ölçüsü.
  \item \term{RAG} --- embedding tabanlı retrieval'ı cevaba bağlar.
  \item \term{Chunking} --- neyin embed edileceğini belirler; vektörün kalitesini etkiler.
\end{itemize}

% ----------------------------------------------------------------------
\stage{12}{Özet}
% ----------------------------------------------------------------------
\begin{itemize}
  \item Embedding, metni \textbf{anlam-koruyan} bir sayı vektörüne çevirir.
  \item Benzer anlam $\approx$ yakın konum; benzerlik bir geometri sorusuna döner.
  \item Yazım yerine \emph{niyeti} yakalar: ``yağmurluk'' $\to$ ``su geçirmez mont''.
  \item İndeks ve sorgu aynı modelle, aynı uzayda üretilmelidir.
  \item Tam terim önemliyse lexical ile birleştir (hybrid).
\end{itemize}

% ----------------------------------------------------------------------
\stage{13}{Mühendislik İlkesi}
% ----------------------------------------------------------------------
\begin{principlebox}
\textbf{İlke \#4 — Anlamı arıyorsan, anlamı temsil eden uzayda ara.}\\[4pt]
\normalsize Harf benzerliği anlam benzerliği değildir. Embedding bu boşluğu
kapatır --- ama yalnızca karşılaştırdığın her vektör \emph{aynı uzaydaysa}.
\end{principlebox}
